% Options for packages loaded elsewhere
\PassOptionsToPackage{unicode}{hyperref}
\PassOptionsToPackage{hyphens}{url}
\PassOptionsToPackage{dvipsnames,svgnames,x11names}{xcolor}
\documentclass[
  11pt,
  a4paper,
]{article}
\usepackage{xcolor}
\usepackage[margin=19mm,headheight=15pt]{geometry}
\usepackage{amsmath,amssymb}
\setcounter{secnumdepth}{-\maxdimen} % remove section numbering
\usepackage{iftex}
\ifPDFTeX
  \usepackage[T1]{fontenc}
  \usepackage[utf8]{inputenc}
  \usepackage{textcomp} % provide euro and other symbols
\else % if luatex or xetex
  \usepackage{unicode-math} % this also loads fontspec
  \defaultfontfeatures{Scale=MatchLowercase}
  \defaultfontfeatures[\rmfamily]{Ligatures=TeX,Scale=1}
\fi
\usepackage{lmodern}
\ifPDFTeX\else
  % xetex/luatex font selection
\fi
% Use upquote if available, for straight quotes in verbatim environments
\IfFileExists{upquote.sty}{\usepackage{upquote}}{}
\IfFileExists{microtype.sty}{% use microtype if available
  \usepackage[]{microtype}
  \UseMicrotypeSet[protrusion]{basicmath} % disable protrusion for tt fonts
}{}
\makeatletter
\@ifundefined{KOMAClassName}{% if non-KOMA class
  \IfFileExists{parskip.sty}{%
    \usepackage{parskip}
  }{% else
    \setlength{\parindent}{0pt}
    \setlength{\parskip}{6pt plus 2pt minus 1pt}}
}{% if KOMA class
  \KOMAoptions{parskip=half}}
\makeatother
\setlength{\emergencystretch}{3em} % prevent overfull lines
\providecommand{\tightlist}{%
  \setlength{\itemsep}{0pt}\setlength{\parskip}{0pt}}

\usepackage{fontspec}
\setmainfont{Liberation Serif}
\setsansfont{Liberation Sans}
\setmonofont{Liberation Mono}[Scale=0.86]
\usepackage{ragged2e,indentfirst,xurl}
\usepackage{longtable,booktabs,array,calc}
\usepackage{xcolor,fancyhdr,titlesec,enumitem}
\definecolor{ink}{HTML}{19354A}
\definecolor{accent}{HTML}{167D8D}
\titleformat{\section}{\Large\sffamily\bfseries\color{ink}}{}{0pt}{}
\titleformat{\subsection}{\normalsize\sffamily\bfseries\color{ink}}{}{0pt}{}
\titleformat{\subsubsection}{\normalsize\sffamily\bfseries\color{ink}}{}{0pt}{}
\titlespacing*{\section}{1.25cm}{0pt}{8pt}
\titlespacing*{\subsection}{1.25cm}{8pt}{4pt}
\titlespacing*{\subsubsection}{1.25cm}{6pt}{3pt}
\setlength{\parindent}{1.25cm}
\setlength{\parskip}{3pt}
\setlength{\emergencystretch}{3em}
\setlength{\tabcolsep}{4pt}
\renewcommand{\arraystretch}{1.08}
\setlist{nosep,leftmargin=1.5em,topsep=3pt}
\pagestyle{fancy}\fancyhf{}
\fancyhead[L]{\small\sffamily ICSI438 / SEMINAR 02}
\fancyhead[R]{\small\sffamily Nogoolin / M2}
\fancyfoot[R]{\small\thepage}
\renewcommand{\headrulewidth}{0.3pt}
\AtBeginDocument{\fontsize{10.5}{12.5}\selectfont\justifying\setlength{\parindent}{1.25cm}}
\usepackage{bookmark}
\IfFileExists{xurl.sty}{\usepackage{xurl}}{} % add URL line breaks if available
\urlstyle{same}
\hypersetup{
  colorlinks=true,
  linkcolor={Maroon},
  filecolor={Maroon},
  citecolor={Blue},
  urlcolor={blue!55!black},
  pdfcreator={LaTeX via pandoc}}

\author{}
\date{}

\newcounter{none}
\begin{document}

\section{SRS Scope Charter}\label{doc1-srs-scope-charter}

\textbf{Nogoolin · ICSI438 / M2 / UE-2 · Тэнгис · 2026-09-21 · v0.2,
review draft}

\subsection{Төслийг сольсон
шийдвэр}\label{doc1-ux442ux4e9ux441ux43bux438ux439ux433-ux441ux43eux43bux44cux441ux43eux43d-ux448ux438ux439ux434ux432ux44dux440}

Би Seminar 1-ийн ажлаа зуны дадлагын Document RAG Chatbot төсөл дээр
хийж, хүлээлгэн өгсөн. Week 2-оос өөрийн бодит бүтээгдэхүүн, portfolio
болох Nogoolin-ийг үндсэн кейсээр сонгов. Энэ төсөлд шаардлага,
архитектур, API, deployment-ийн баримтыг одоогийн кодтой нийцүүлэх
хэрэгцээ байгаа тул M4, M5, M8, M9-ийн ажлыг бодит төсөлдөө үргэлжлүүлэх
боломжтой. Seminar 1-ийг дахин хийхгүй; шилжилтийн persona холбоог
M2-ийн нэмэлтээр ил тод тэмдэглэнэ.

\textbf{Public repository:}
\href{https://github.com/Tengis01/Nogoolin}{github.com/Tengis01/Nogoolin}

\subsection{Зорилго ба
уншигч}\label{doc1-ux437ux43eux440ux438ux43bux433ux43e-ux431ux430-ux443ux43dux448ux438ux433ux447}

Каталогийн хэрэглэгч бүтээгдэхүүний тухай асуулга илгээх, хадгалсан
бүтээгдэхүүнээ дахин харах боломжийг тодорхойлно. Гол уншигч нь
хэрэгжүүлж, шалгах хөгжүүлэгч Тэнгис; хоёрдогч уншигч нь хичээлийн багш.
Phase 0 баримтыг эх материал болгон ашиглаж, технологийн сонголтыг
SDD-д, шалгаж болох үр дүнг SRS-д тусгана.

\subsection{Included --- энэ
хүрээнд}\label{doc1-included-ux44dux43dux44d-ux445ux4afux440ux44dux44dux43dux434}

\begin{itemize}
\tightlist
\item
  \textbf{Inquiry + wishlist:} guest болон нэвтэрсэн хэрэглэгчийн
  асуулга, хүлээн авсан баталгаа; нэвтэрсэн хэрэглэгчийн
  хадгалах/харах/хасах урсгал. Админ inbox нь асуулга очсоныг шалгах
  төгсгөлийн цэг.
\item
  \textbf{Өгөгдлийн тусгаарлалт:} хэрэглэгч өөрийн wishlist, inquiry
  history-г л уншиж/өөрчилнө; админ эрхийг тусдаа шалгана.
\item
  \textbf{Hero fallback:} 3D asset бэлэн бус үед орлуулагч, WebGL-гүй
  болон хөдөлгөөн багасгах тохиргоонд статик хувилбар; каталогт хүрэх
  боломж.
\item
  \textbf{Шалгах ба build хийх хүрээ:} API-ийн local Docker build,
  deploy-оос өмнөх тестийн зайлшгүй хаалт. Docker-ийн сонголт SDD-д,
  амжилт/алдааны шалгуур NFR-03-д байна.
\end{itemize}

\subsection{Excluded --- M2/M3-ийн энэ багцаас
хассан}\label{doc1-excluded-m2m3-ux438ux439ux43d-ux44dux43dux44d-ux431ux430ux433ux446ux430ux430ux441-ux445ux430ux441ux441ux430ux43d}

\textbf{Бүтээгдэхүүн бүрийн 360° viewer:} тус бүрд нь 3D asset бэлтгэх
ажил ганцаараа хөгжүүлэх MVP-д хэт их тул энгийн зургийн gallery
ашиглана (D-07). Энэ нь нүүрний 3D intro-г хассан шийдвэр биш.
Viewer-ийг post-MVP-д дахин авч үзнэ; энэ багцад хэрэгжүүлэх, тестлэх
үүрэг үүсгэхгүй.

\subsection{Postponed --- дараа
хийх}\label{doc1-postponed-ux434ux430ux440ux430ux430-ux445ux438ux439ux445}

\textbf{Phase 5:} захиалга, төлбөр/checkout, хүргэлтийн ажиллагаа; одоо
\nolinkurl{delivery_enabled=false} гэсэн шийдвэр хүчинтэй. \textbf{Phase
6:} cloud Supabase холбох, production deployment баталгаажуулах, App
Store/Play Store нийтлэх. Mobile inquiry/wishlist, бодит GLB/.riv asset
болон бүх каталог/admin шаардлагыг энэ таван шаардлага бүрэн хамрахгүй;
бүтээгдэхүүний backlog-д үлдээнэ.

\subsection{Хүрээ тогтоох үндэслэл ба дуусах
нөхцөл}\label{doc1-ux445ux4afux440ux44dux44d-ux442ux43eux433ux442ux43eux43eux445-ux4afux43dux434ux44dux441ux43bux44dux43b-ux431ux430-ux434ux443ux443ux441ux430ux445-ux43dux4e9ux445ux446ux4e9ux43b}

Chinchilla-ийн х. 53--55 дахь аргаас: \textbf{хэн юу хүлээж байна?} ---
хэрэглэгч баталгаа, хөгжүүлэгч давтаж шалгах нөхцөл; \textbf{ямар далд
таамаг байна?} --- нэвтэрсэн л бол бүх өгөгдөл харах эрхтэй биш, asset
болон cloud бэлэн гэж үзэхгүй; \textbf{эхлээд юуг баримтжуулах вэ?} ---
эхлэлээс үр дүн хүртэлх inquiry/wishlist урсгал. Эдгээр нь номын аргыг
төсөлдөө хэрэглэсэн асуултууд, үгчилсэн эшлэл биш.

M2-ийн бэлэн багц: SRS/SDD загвар, 5 шаардлага, 8 баганат traceability,
persona холбоо, review draft. Ярилцлага хийгдээгүй; peer review илгээх,
нийтлэх/submission алхмыг хэрэглэгчийн заавраар алгассан. W1 persona-ийн
шинэ төсөлд шилжүүлсэн холбоог reviewer-ээр баталгаажуулах шаардлагатай.

\textbf{Эх:} \href{../../phase-0/01-vision.md}{Vision},
\href{../../phase-0/10-roadmap.md}{Roadmap},
\href{../../../agent-context/DECISIONS.md}{D-07 / ADR-011},
\href{Software_Project_Documentation_Week_02_Handout.pdf}{W2 handout};
Chinchilla, бүлэг 5, хэвлэмэл х. 53--55 (локал номын PDF х. 74--76).

\clearpage
\section{Nogoolin persona --- Week 2
нэмэлт}\label{doc2-nogoolin-persona-week-2-ux43dux44dux43cux44dux43bux442}

\textbf{Тэнгис · 2026-09-15 · US-2.3 · Өөрийн ажиглалт.}

\subsection{Уншигч ба бүтээгдэхүүний
хэрэглэгч}\label{doc2-ux443ux43dux448ux438ux433ux447-ux431ux430-ux431ux4afux442ux44dux44dux433ux434ux44dux445ux4afux4afux43dux438ux439-ux445ux44dux440ux44dux433ux43bux44dux433ux447}

\textbf{Хөгжүүлэгч persona:} Тэнгис, TypeScript/React ашигладаг,
web/API/өгөгдлийн санг холбож ажиллуулна. Зорилго нь Nogoolin-ийн
inquiry/wishlist урсгалыг local орчинд давтаж шалгаж, алдааг аль
давхаргаас хайхаа ойлгох. Хэрэгсэл: Fedora, VS Code, Git, pnpm, Docker,
DevTools.

\textbf{Бүтээгдэхүүний хэрэглэгч:} каталог үзэж, асуулга илгээх эсвэл
дараа үзэх бүтээгдэхүүнээ хадгалах хүн. Энэ нь
\href{../../phase-0/01-vision.md}{Phase 0 vision}-ийн catalog
visitor/inquiry customer тайлбарын хураангуй; шинэ хүнтэй ярилцсан үр
дүн биш.

\subsection{P-NG-01 --- Session ба өгөгдөл эзэмших эрхийг
ялгах}\label{doc2-p-ng-01-session-ux431ux430-ux4e9ux433ux4e9ux433ux434ux4e9ux43b-ux44dux437ux44dux43cux448ux438ux445-ux44dux440ux445ux438ux439ux433-ux44fux43bux433ux430ux445}

Нэвтрэлт ажиллаж байсан ч A хэрэглэгч B-ийн wishlist, inquiry history-г
харах ёсгүй. Шалгах зааварт guest, customer A/B, admin-ийн ялгаа байхгүй
бол эрхийн алдаа анзаарагдахгүй үлдэнэ. Энэ бол код болон security
баримтаас гаргасан эрсдэлийн ажиглалт; бодит халдлага болсон гэсэн үг
биш.

\textbf{Source:}
\href{../../phase-0/08-security.md\#9-idor-prevention}{08-security §9},
\href{../../../backend/api/src/repositories/supabase/cart.repository.ts}{cart
repository},
\href{../../../backend/api/src/repositories/supabase/inquiry.repository.ts}{inquiry
repository}. \textbf{Холбоо:}
\href{../../requirements/requirements.md\#nfr-01}{NFR-01}.

\subsection{P-NG-02 --- Asset бэлэн бус үед урсгалыг
шалгах}\label{doc2-p-ng-02-asset-ux431ux44dux43bux44dux43d-ux431ux443ux441-ux4afux435ux434-ux443ux440ux441ux433ux430ux43bux44bux433-ux448ux430ux43bux433ux430ux445}

Бодит Green Tara GLB/.riv бэлэн болоогүй үед каталогийн урсгалыг
хүлээлгэхгүй шалгах хэрэгтэй. Asset тохируулаагүй нөхцөл дэх орлуулагч,
WebGL-гүй болон reduced-motion хувилбарын үр дүнг тус бүр тодорхой
болгоно.

\textbf{Source:} \href{../../../agent-context/ASSET_SPECS.md}{Asset
contract},
\href{../../../apps/web/src/components/intro/deity.tsx}{deity.tsx},
\href{../../../apps/web/src/components/intro/hero-intro.tsx}{hero-intro.tsx}.
\textbf{Холбоо:}
\href{../../requirements/requirements.md\#nfr-02}{NFR-02}.

\subsection{P-NG-03 --- Хуучин заавар ба бодит орчин
зөрөх}\label{doc2-p-ng-03-ux445ux443ux443ux447ux438ux43d-ux437ux430ux430ux432ux430ux440-ux431ux430-ux431ux43eux434ux438ux442-ux43eux440ux447ux438ux43d-ux437ux4e9ux440ux4e9ux445}

Flutter/Express гэсэн хуучин тайлбаруудыг одоогийн React Native/Fastify
шийдвэртэй андуурч болно. Мөн local өгөгдлийн сан байхгүй үед тест
алгасагдсаныг бүрэн шалгалт өнгөрсөн гэж ойлгох эрсдэлтэй.

\textbf{Source:}
\href{../../../agent-context/DECISIONS.md}{ADR-003/004/011},
\href{../../../backend/api/tests/inquiry-cart.integration.test.ts}{integration
test-ийн skip салаа}. \textbf{Холбоо:}
\href{../../requirements/requirements.md\#nfr-03}{NFR-03},
\href{../../architecture/sdd-template.md}{SDD}.

\subsection{Week 1-тэй холбосон
үндэслэл}\label{doc2-week-1-ux442ux44dux439-ux445ux43eux43bux431ux43eux441ux43eux43d-ux4afux43dux434ux44dux441ux43bux44dux43b}

\href{../sem1/wiki/02_audience_persona.md}{Өмнө хүлээлгэн өгсөн W1
persona}-ийн \textbf{P1} нь нэвтрэлтийн cookie/session-ийг шалгах
зааврын хэрэгцээ байсан. Түүний authentication-ийг оношлох хэрэгцээг
Nogoolin-д өргөжүүлж \textbf{P1 → P-NG-01 → NFR-01} гэж холбов.
Authentication болон ownership нь өөр ойлголт; W1-д IDOR байсан гэж
үзээгүй. Мөн W1-ийн \textbf{P3} нь local ба pipeline-д өөр үр дүн
гарахыг оношлох, өөр замтай байх хэрэгцээ; үүнээс \textbf{P3 → P-NG-03 →
NFR-03} холбоог шинэ төсөлд шилжүүлэн гаргав.

Эдгээр нь Week 2-т нэмсэн тайлбарласан холбоо. US-2.3-ийн W1 pain point
шалгуурт энэ шилжүүлэлт нийцэх эсэхийг багш/reviewer-ээр баталгаажуулна;
Week 1-ийн эхийг өөрчлөөгүй.

\clearpage
\section{SRS загвар -
Nogoolin}\label{doc3-srs-ux437ux430ux433ux432ux430ux440---nogoolin}

\textbf{M2 / US-2.1 · Тэнгис · 2026-09-21 · v0.2 · SRS загвар}

\subsection{Загварыг ашиглах
нь}\label{doc3-ux437ux430ux433ux432ux430ux440ux44bux433-ux430ux448ux438ux433ux43bux430ux445-ux43dux44c}

IEEE 830-1998-ийн агуулгын бүтцийг ISO/IEC/IEEE 29148:2018-ийн шаардлага
тодорхойлох, шалгах, мөрдөх зарчимтай хослуулна. Доорх
\emph{{[}\ldots{]}} хэсэгт M3-аас агуулга нэмнэ. M2-д таван жишээ
шаардлагыг тусад нь бичсэн; энэ нь бүрэн SRS биш.

\textbf{Баримтын бүртгэл:} төсөл Nogoolin; зохиогч Тэнгис; хувилбар
v0.2; төлөв draft; хянагч \emph{{[}нэр{]}}; баталсан огноо
\emph{{[}огноо{]}}.

\subsubsection{1. Оршил}\label{doc3-ux43eux440ux448ux438ux43b}

\textbf{1.1 Зорилго.} \emph{{[}Энэ SRS ямар шийдвэр, хэрэгжилт,
шалгалтад ашиглагдах вэ?{]}} Гол уншигч нь хөгжүүлэгч Тэнгис; хоёрдогч
уншигч нь хичээлийн багш.

\textbf{1.2 Хамрах хүрээ.}
\href{../../requirements/scope-charter.md}{Scope Charter}-ийг мөрдөнө.
Эхлэх хүрээ: inquiry, wishlist, хувийн өгөгдлийн тусгаарлалт, hero
fallback, API тестийн хаалт. \emph{{[}Системийн хил ба амжилтын
хэмжүүр.{]}}

\textbf{1.3 Нэр томьёо.} \emph{{[}Нэр - нэг утгатай тайлбар.{]}} Эхлэл:
inquiry = бүтээгдэхүүний асуулга; wishlist = дараа үзэхээр хадгалсан
жагсаалт; FR = үйлдлийн шаардлага; NFR = чанарын шаардлага.

\textbf{1.4 Эх сурвалж.} \emph{{[}Баримтын нэр, хувилбар/огноо, хэсэг,
холбоос.{]}} Эхлэл: \href{https://github.com/Tengis01/Nogoolin}{Nogoolin
public repository}, W2 handout, Lecture 02,
\href{../../phase-0/README.md}{Phase 0 эхүүд}, IEEE 830, ISO/IEC/IEEE
29148.

\textbf{1.5 Бүтцийн тайлбар.} §2-т орчин ба хэрэглэгч, §3-т шаардлага;
хавсралтуудад mapping, чанарын checklist, traceability байна.
Технологийн шийдвэрийг \href{../../architecture/sdd-template.md}{SDD
загварт} бөглөнө.

\subsubsection{2. Ерөнхий
тодорхойлолт}\label{doc3-ux435ux440ux4e9ux43dux445ux438ux439-ux442ux43eux434ux43eux440ux445ux43eux439ux43bux43eux43bux442}

\textbf{2.1 Бүтээгдэхүүний орчин.} \emph{{[}Системийн зориулалт, хил,
гаднын системтэй харилцах байдал.{]}} Эхлэл: Nogoolin нь шашин, зан
үйлийн бүтээгдэхүүний каталог.

\textbf{2.2 Гол үйлдлүүд.} \emph{{[}Хэрэглэгчийн зорилгоор бүлэглэсэн
үйлдлийн жагсаалт; FR ID холбоос.{]}} Эхлэх жишээ: асуулга илгээх FR-01,
бүтээгдэхүүн хадгалах FR-02.

\textbf{2.3 Хэрэглэгчийн онцлог.} \emph{{[}Дүр, зорилго, мэдлэг,
хүндрэл, W1 persona-ийн эх.{]}} Guest, customer, admin-ийг ялгана;
\href{wiki/persona-and-pivot.md}{persona нэмэлт}-ээс эхэлнэ.

\textbf{2.4 Хязгаарлалт.} \emph{{[}Бизнес, орчин, нөөцийн зайлшгүй
хязгаар ба үндэслэл.{]}} Эхлэл: ганцаар хөгжүүлэх, монгол интерфэйс,
хүргэлт идэвхгүй.

\textbf{2.5 Таамаг ба хамаарал.} \emph{{[}Таамаг, эзэн, шалгах хугацаа,
буруу бол нөлөө.{]}} Local DB болон тестийн өгөгдөл бэлэн эсэхийг
шалгана; cloud, asset, stakeholder approval-ийг бэлэн гэж таамаглахгүй.

\textbf{2.6 Дараа хийх хүрээ.} \emph{{[}Боломж, хойшлуулах шалтгаан,
зорилтот үе.{]}} Scope Charter-ийн postponed жагсаалтыг холбоно.
Хичээлийн W3 ба бүтээгдэхүүний Phase 3-ыг адилтгахгүй.

\subsubsection{3. Тодорхой
шаардлагууд}\label{doc3-ux442ux43eux434ux43eux440ux445ux43eux439-ux448ux430ux430ux440ux434ux43bux430ux433ux443ux443ux434}

\textbf{3.1 Гадаад интерфэйс.} \emph{{[}Оролт/гаралт, формат, алдаа,
харилцах тал.{]}} 3.1.1 хэрэглэгчийн, 3.1.2 төхөөрөмжийн, 3.1.3
программын, 3.1.4 холбооны интерфэйс гэж салгана. Төхөөрөмжийн
интерфэйсийн n-a үндэслэлийг mapping-д тэмдэглэв.

\textbf{3.2 Функциональ шаардлага.} \emph{{[}Доорх шаардлагын картаар FR
бүрийг бич.{]}} M2-ийн \href{../../requirements/requirements.md}{FR-01,
FR-02} нь жишээ; үлдсэнийг W3-д нэмнэ.

\textbf{3.3 Гүйцэтгэл.} \emph{{[}Ачаалал, өгөгдлийн хэмжээ, орчин,
хугацааны хэмжүүр, pass/fail босго.{]}} ``Хурдан'' гэх мэт хэмжүүргүй
өгүүлбэр хэрэглэхгүй.

\textbf{3.4 Логик өгөгдөл.} \emph{{[}Хадгалах мэдээлэл, холбоо,
unique/ownership дүрэм, хадгалах хугацаа.{]}} Өгөгдлийн сангийн
бүтээгдэхүүн сонгох шийдвэрийг SDD-д бичнэ.

\textbf{3.5 Дизайнд тавих зайлшгүй хязгаар.} \emph{{[}Гаднаас тогтоосон
техникийн хязгаар, түүний эх ба үндэслэл.{]}} Өөрийн сонгосон
хэрэгжүүлэлтийн аргыг энд шаардлага болгож хуулж болохгүй.

\textbf{3.6 Чанарын шаардлага.} \emph{{[}Найдвартай ажиллагаа, бэлэн
байдал, аюулгүй байдал, засварлах боломж, зөөврийн байдал тус бүрийн
NFR.{]}} M2-ийн NFR-01\ldots03-ыг холбож, дутагдах хэсгийг W3-д нөхнө.

\textbf{3.7 Бусад шаардлага.} \emph{{[}Хувийн мэдээлэл, audit, хэрэглэх
бодлогын эх ба шалгуур.{]}} Мэдээлэл дутуу хэсгийг planned гэж
тэмдэглэнэ; тодорхойгүйг n-a гэж үзэхгүй.

\subsection{Хавсралт A. IEEE 830 section
mapping}\label{doc3-ux445ux430ux432ux441ux440ux430ux43bux442-a.-ieee-830-section-mapping}

\textbf{Төлөв:} \nolinkurl{complete} = M2-т шаардсан тухайн агуулга
бэлэн; бүтээгдэхүүн хэрэгжсэн гэсэн үг биш. \nolinkurl{planned (W…)} =
агуулгыг нөхөх үе. \nolinkurl{n-a (justified)} = хүрээнд хамаарахгүй,
нэг өгүүлбэр үндэслэлтэй. Мөр бүр яг нэг төлөвтэй.

\begin{longtable}[]{@{}
  >{\raggedright\arraybackslash}p{(\linewidth - 4\tabcolsep) * \real{0.2600}}
  >{\raggedright\arraybackslash}p{(\linewidth - 4\tabcolsep) * \real{0.5500}}
  >{\raggedright\arraybackslash}p{(\linewidth - 4\tabcolsep) * \real{0.1900}}@{}}
\toprule\noalign{}
\begin{minipage}[b]{\linewidth}\raggedright
Section
\end{minipage} & \begin{minipage}[b]{\linewidth}\raggedright
Content / агуулга, эх
\end{minipage} & \begin{minipage}[b]{\linewidth}\raggedright
Status
\end{minipage} \\
\midrule\noalign{}
\endhead
\bottomrule\noalign{}
\endlastfoot
1.1 Purpose & §1.1-ийн зорилго, уншигчийн хэрэгцээг эцэслэх & planned
(W03) \\
1.2 Scope & Нэг хуудас Scope Charter & complete \\
1.3 Definitions & Нэр томьёоны нэгдсэн тайлбар & planned (W03) \\
1.4 References & W2/лекц/стандарт/Phase 0 эхийн холбоос & complete \\
1.5 Overview & §1.5-ийн бүтцийн тайлбар & complete \\
2.1 Product perspective & Системийн хил, гаднын орчин; Vision & planned
(W03) \\
2.2 Product functions & FR-01/02-оос үйлдлийн хүрээг өргөжүүлэх &
planned (W03) \\
2.3 User characteristics & Persona нэмэлт, ярилцлагаар нягтлах & planned
(W03) \\
2.4 Constraints & Scope Charter, нөөц ба бизнесийн хязгаар & planned
(W03) \\
2.5 Assumptions/dependencies & Таамаг бүрийн эзэн, шалгах нөхцөл &
planned (W03) \\
2.6 Apportioning & Scope Charter-ийн postponed жагсаалт & complete \\
3.1.1 User interfaces & Хэрэглэгчийн урсгал, оролт, баталгаа/алдаа &
planned (W03) \\
3.1.2 Hardware interfaces & Энэ web/API хүрээнд тусгай төхөөрөмж,
сенсортой холболт байхгүй. & n-a (justified) \\
3.1.3 Software interfaces & Гаднын систем ба интерфэйсийн шаардлага;
нарийн API contract W5 & planned (W03) \\
3.1.4 Communications & Протокол, өгөгдлийн формат, хамгаалалтын
шаардлага & planned (W03) \\
3.2 Functional requirements & Хоёр FR жишээ; W3-д 15 FR болгох & planned
(W03) \\
3.3 Performance & Ачаалал, орчин, хэмжүүр, босго & planned (W03) \\
3.4 Logical database & Өгөгдөл ба бүрэн бүтэн байдлын дүрэм & planned
(W03) \\
3.5 Design constraints & Зайлшгүй хязгаар ба эх сурвалж & planned
(W03) \\
3.6.1 Reliability & Алдаа/сэргэлтийн шалгуур; NFR-02 эхлэл & planned
(W03) \\
3.6.2 Availability & Хэмжих хугацаа, бэлэн байдлын зорилт & planned
(W03) \\
3.6.3 Security & NFR-01; эрх ба хувийн өгөгдөл & planned (W03) \\
3.6.4 Maintainability & NFR-03; шаардлагыг W3, CI нотолгоог W9 & planned
(W03) \\
3.6.5 Portability & Дэмжих орчин ба шалгах нөхцөл & planned (W03) \\
3.7 Other requirements & Хувийн мэдээлэл, audit бодлого & planned
(W03) \\
Appendix B & Шаардлагын карт ба найман чанарын шалгуур & complete \\
Appendix C & 5 мөртэй, 8 баганат traceability; W1 холбоо & complete \\
\end{longtable}

\subsection{Хавсралт B. Шаардлагын карт ба чанарын
шалгалт}\label{doc3-ux445ux430ux432ux441ux440ux430ux43bux442-b.-ux448ux430ux430ux440ux434ux43bux430ux433ux44bux43d-ux43aux430ux440ux442-ux431ux430-ux447ux430ux43dux430ux440ux44bux43d-ux448ux430ux43bux433ux430ux43bux442}

Энэ картыг §3-ын шаардлага бүрд хуулж бөглөнө. ISO/IEC/IEEE 29148-ийн
шаардлагын шинж, атрибут, удирдлагын зарчмыг хэрэгжүүлэхэд ашиглана;
доорх бүх талбар стандартын үгчилсэн хуулбар биш.

\begin{longtable}[]{@{}
  >{\raggedright\arraybackslash}p{(\linewidth - 2\tabcolsep) * \real{0.2500}}
  >{\raggedright\arraybackslash}p{(\linewidth - 2\tabcolsep) * \real{0.7500}}@{}}
\toprule\noalign{}
\begin{minipage}[b]{\linewidth}\raggedright
Талбар
\end{minipage} & \begin{minipage}[b]{\linewidth}\raggedright
Загвар
\end{minipage} \\
\midrule\noalign{}
\endhead
\bottomrule\noalign{}
\endlastfoot
ID / төрөл / нэр & \emph{{[}FR-XX эсвэл NFR-XX{]} / {[}FR эсвэл NFR{]} /
{[}нэг үр дүнгийн нэр{]}} \\
Шаардлагын өгүүлбэр & \emph{{[}Нөхцөл{]} үед систем {[}ажиглагдах нэг үр
дүн{]}-г {[}must/shall/may{]} хангана.} \\
Priority & \emph{{[}must / shall / may{]}}; хичээлийн тэмдэглэгээг
мөрдөж \nolinkurl{should} хэрэглэхгүй \\
Source / үндэслэл & \emph{{[}Баримт/хэсэг/холбоос, persona pain point,
яагаад хэрэгтэй{]}} \\
Owner / source interview & \emph{{[}Нэр{]} / {[}ярилцсан хүн, огноо,
тэмдэглэлийн холбоос; хийгдээгүй бол хийгдээгүй{]}} \\
Acceptance criteria & \emph{Given {[}урьдчилсан нөхцөл{]}, When
{[}үйлдэл{]}, Then {[}ажиглах үр дүн ба pass/fail босго{]}.} \\
Verification & \emph{{[}Test / Inspection / Analysis{]} + {[}шалгалтын
ID, алхам, орчин, нотолгоо{]}} \\
Dependency / Risk & \emph{{[}Хамаарах ID/нөхцөл{]} / {[}эрсдэл{]}} \\
Status / Last-Reviewed & \emph{{[}Draft / reviewed / approved \ldots{]}
/ {[}YYYY-MM-DD{]}} \\
\end{longtable}

\subsubsection{W2-ын найман чанарын
шалгуур}\label{doc3-w2-ux44bux43d-ux43dux430ux439ux43cux430ux43d-ux447ux430ux43dux430ux440ux44bux43d-ux448ux430ux43bux433ux443ux443ux440}

\begin{longtable}[]{@{}
  >{\raggedright\arraybackslash}p{(\linewidth - 2\tabcolsep) * \real{0.2500}}
  >{\raggedright\arraybackslash}p{(\linewidth - 2\tabcolsep) * \real{0.7500}}@{}}
\toprule\noalign{}
\begin{minipage}[b]{\linewidth}\raggedright
Шалгуур
\end{minipage} & \begin{minipage}[b]{\linewidth}\raggedright
Хянах асуулт
\end{minipage} \\
\midrule\noalign{}
\endhead
\bottomrule\noalign{}
\endlastfoot
Correct & Хэрэгцээ нь эх сурвалжаар нотлогдож, stakeholder-аар
нягтлагдсан уу? \\
Unambiguous & Нэр томьёо, нөхцөл, үр дүнг нэг утгаар ойлгох уу? \\
Complete & Тухайн шаардлагын урьдчилсан нөхцөл, хэвийн/алдааны үр дүн
хангалттай юу? \\
Consistent & Бусад ID, Scope Charter-тай зөрчилдөх үү? \\
Ranked & Priority ил тод уу; өөрчлөгдөх магадлалыг шаардлагатай бол
тэмдэглэсэн үү? \\
Verifiable & Test/Inspection/Analysis-аар pass/fail-ийг тогтоож болох
уу? \\
Modifiable & Нэг ID нэг үүрэгтэй, давхар хуулбаргүй, засварлахад
ойлгомжтой юу? \\
Traceable & Эх хэрэгцээ → шаардлага → шалгалтын холбоо байгаа юу? \\
\end{longtable}

Энэ найман нэрийг Lecture 02-ын rubric-аар хэрэглэв. IEEE 830 §4.3-т мөн
эдгээр SRS чанарыг нэрлэдэг; ISO 29148-ийн бүх шинжийг зөвхөн энэ
наймаар хязгаарлаж ойлгохгүй.

\subsection{Хавсралт C. Traceability ба
өөрчлөлт}\label{doc3-ux445ux430ux432ux441ux440ux430ux43bux442-c.-traceability-ux431ux430-ux4e9ux4e9ux440ux447ux43bux4e9ux43bux442}

\href{../../requirements/traceability-matrix.md}{Traceability
Matrix}-ийн баганыг
\nolinkurl{ID | Source | Owner | Verification | Dependency | Risk | Status | Last-Reviewed}
гэсэн дарааллаар хадгална. Шаардлага бүр нэг мөртэй, дор хаяж гурван
багана бөглөсөн байна. M2-ийн таван жишээнд найман баганыг бөглөсөн.
Нэгээс доошгүй мөр W1 persona pain point-той холбоотой байна; төсөл
сольсон холбоог тайлбарлана.

Өөрчлөхдөө ID-г хадгалж, source, шалгуур, matrix, review огноог хамтад
нь шинэчилнэ. Тестийн нотолгоо байхгүй бол verified гэж тэмдэглэхгүй.
Бүрэн SRS (15 FR + 5 NFR) нь W3/M3-ын ажил.

\textbf{Эх:}
\href{Software_Project_Documentation_Week_02_Handout.pdf}{W2 handout},
\href{../lecture/Lecture_02\%20Why\%20Documentation\%20Documentation\%20Standards.pdf}{Lecture
02, слайд 4--8};
\href{https://www.math.uaa.alaska.edu/~afkjm/cs401/IEEE830.pdf}{IEEE
830-1998}, §4.3, §5;
\href{https://standards.ieee.org/ieee/29148/6937/}{ISO/IEC/IEEE
29148:2018}. Хичээлийн унших чиглэл: Bhatti бүлэг 2, х. 23--44 (хүрээ),
бүлэг 4, х. 67--81 (шаардлагын эх ба уншигч); Chinchilla бүлэг 5, х.
53--55 (хүлээлт, таамаг, одоогийн дуусах нөхцөл).

\clearpage
\begingroup\footnotesize\setlength{\tabcolsep}{2pt}
\section{Traceability Matrix}\label{doc5-traceability-matrix}

\textbf{Nogoolin · M2 / US-2.3 · v0.2 · 2026-09-21}

Доорх нь яг 8 баганатай, 5 мөртэй. Last-Reviewed нь баримт/кодыг уншсан
огноо; тест ажиллуулсан огноо биш. Owner нь Тэнгис (хөгжүүлэгч, баримт
бичигч). Бүх шаардлага draft, stakeholder interview хийгдээгүй.

\begin{longtable}[]{@{}
  >{\raggedright\arraybackslash}p{(\linewidth - 14\tabcolsep) * \real{0.0650}}
  >{\raggedright\arraybackslash}p{(\linewidth - 14\tabcolsep) * \real{0.2550}}
  >{\raggedright\arraybackslash}p{(\linewidth - 14\tabcolsep) * \real{0.0650}}
  >{\raggedright\arraybackslash}p{(\linewidth - 14\tabcolsep) * \real{0.1200}}
  >{\raggedright\arraybackslash}p{(\linewidth - 14\tabcolsep) * \real{0.1150}}
  >{\raggedright\arraybackslash}p{(\linewidth - 14\tabcolsep) * \real{0.1450}}
  >{\raggedright\arraybackslash}p{(\linewidth - 14\tabcolsep) * \real{0.1400}}
  >{\raggedright\arraybackslash}p{(\linewidth - 14\tabcolsep) * \real{0.0950}}@{}}
\toprule\noalign{}
\begin{minipage}[b]{\linewidth}\raggedright
ID
\end{minipage} & \begin{minipage}[b]{\linewidth}\raggedright
Source
\end{minipage} & \begin{minipage}[b]{\linewidth}\raggedright
Owner
\end{minipage} & \begin{minipage}[b]{\linewidth}\raggedright
Verification
\end{minipage} & \begin{minipage}[b]{\linewidth}\raggedright
Dependency
\end{minipage} & \begin{minipage}[b]{\linewidth}\raggedright
Risk
\end{minipage} & \begin{minipage}[b]{\linewidth}\raggedright
Status
\end{minipage} & \begin{minipage}[b]{\linewidth}\raggedright
Last-Reviewed
\end{minipage} \\
\midrule\noalign{}
\endhead
\bottomrule\noalign{}
\endlastfoot
\href{../../requirements/requirements.md\#fr-01}{FR-01} &
\href{../../phase-0/02-requirements.md}{FR-INQ-001\ldots003};
\href{../../../packages/validation-schemas/src/inquiry.schema.ts}{schema};
\href{../../../backend/api/src/controllers/inquiry.controller.ts}{controller}
& Тэнгис & Test: T-01;
\href{../../../backend/api/tests/inquiry-cart.integration.test.ts}{тестийн
эх} & Бүтээгдэхүүн, local DB, validation, rate limit & Email/message-ийн
хуучин баримтын зөрүү; давхардсан илгээлт & Draft; runtime pending &
2026-09-21 \\
\href{../../requirements/requirements.md\#fr-02}{FR-02} &
\href{../../../backend/api/src/controllers/cart.controller.ts}{cart
controller};
\href{../../../supabase/migrations/20260726120000_inquiry_customer_and_cart.sql}{migration}
& Тэнгис & Test: T-02;
\href{../../../backend/api/tests/inquiry-cart.integration.test.ts}{тестийн
эх} & Auth, published product, unique pair, NFR-01 & Wishlist-ийг
checkout-тай андуурах & Draft; runtime pending & 2026-09-21 \\
\href{../../requirements/requirements.md\#nfr-01}{NFR-01} &
\href{../sem1/wiki/02_audience_persona.md}{W1 persona P1} →
\href{wiki/persona-and-pivot.md\#p-ng-01--session-ба-өгөгдөл-эзэмших-эрхийг-ялгах}{P-NG-01};
\href{../../phase-0/08-security.md\#9-idor-prevention}{08 §9};
\href{../../../backend/api/src/repositories/supabase/cart.repository.ts}{cart
repo};
\href{../../../backend/api/src/repositories/supabase/inquiry.repository.ts}{inquiry
repo} & Тэнгис & Inspection: I-01; A/B runtime нэмэлт & Session
identity, ownership filter, RLS, RBAC & Өөр хүний мэдээлэл алдагдах; W1
холбооны зөвшөөрөл pending & Draft; partial inspection & 2026-09-21 \\
\href{../../requirements/requirements.md\#nfr-02}{NFR-02} &
\href{wiki/persona-and-pivot.md\#p-ng-02--asset-бэлэн-бус-үед-урсгалыг-шалгах}{P-NG-02};
\href{../../phase-0/07-uiux-wireframes.md}{wireframes};
\href{../../../apps/web/src/components/intro/hero-intro.tsx}{hero} &
Тэнгис & Test: T-03; browser evidence pending & Hero, fallback, catalog
route & Asset байхгүйгээс хэрэглэгч каталогт хүрэхгүй & Draft; runtime
pending & 2026-09-21 \\
\href{../../requirements/requirements.md\#nfr-03}{NFR-03} &
\href{../sem1/wiki/02_audience_persona.md}{W1 P3} →
\href{wiki/persona-and-pivot.md\#p-ng-03--хуучин-заавар-ба-бодит-орчин-зөрөх}{P-NG-03};
\href{../../phase-0/09-deployment.md}{09 deployment};
\href{../../../.github/workflows/api-deploy.yml}{workflow};
\href{../../../backend/api/tests/inquiry-cart.integration.test.ts}{skip
branch} & Тэнгис & Inspection: I-02; CI negative-case logs pending &
Local DB in CI, test exit status, Docker build & Бүх тест skip боловч
deploy eligible болох & Gap; planned (W09) & 2026-09-21 \\
\end{longtable}

\subsection{Эх сурвалжийн үнэн зөв
байдал}\label{doc5-ux44dux445-ux441ux443ux440ux432ux430ux43bux436ux438ux439ux43d-ux4afux43dux44dux43d-ux437ux4e9ux432-ux431ux430ux439ux434ux430ux43b}

W1-ийн P1 нь cookie/session асуудал бөгөөд IDOR биш. NFR-01 рүү холбосон
нь Week 2-т нэмсэн, тайлбарласан шинэ хэрэглээ. W1 P3-ийн local/pipeline
ялгаанаас NFR-03-ийн тестийн орчны эрсдэлийг гаргасан. Багш эдгээр
холбоог зөвшөөрсөн гэж үзээгүй;
\href{wiki/peer-review-draft.md}{peer-review draft}-д шалгуулах асуулт
бий. Ярилцлагын source зохиогоогүй.

\subsection{Шинэчлэх
дүрэм}\label{doc5-ux448ux438ux43dux44dux447ux43bux44dux445-ux434ux4afux440ux44dux43c}

Шаардлага өөрчлөгдвөл ID-г хадгалж, эх холбоос, verification-ийн
pass/fail, dependency, risk, status, огноог хамтад нь шинэчилнэ. Тестийн
run URL/log нэмсний дараа л verification status-ийг verified болгоно.
Confluence-ийн Requirements хэсэгт энэ файлын GitHub permalink-ийг
оруулахад бэлэн; хэрэглэгчийн заавраар нийтлэлт/submission-ийг алгассан,
remote commit хийгээгүй.

\endgroup
\clearpage
\section{SDD загвар -
Nogoolin}\label{doc6-sdd-ux437ux430ux433ux432ux430ux440---nogoolin}

\textbf{M2 · Тэнгис · 2026-09-21 · v0.2 · Бөглөх загвар}

SRS нь системээс шаардах үр дүнг, SDD нь түүнийг хэрэгжүүлэх аргыг
тодорхойлно. Lecture 02-ын M2 заалтын дагуу энэ бүтэц бэлтгэв. Доорх
асуултуудыг дараагийн ажлаар бөглөнө; бүрэн архитектурын баримт биш.

\subsection{1. Зорилго ба хамрах
хүрээ}\label{doc6-ux437ux43eux440ux438ux43bux433ux43e-ux431ux430-ux445ux430ux43cux440ux430ux445-ux445ux4afux440ux44dux44d}

\emph{{[}Бөглөх: ямар SRS хувилбар, FR/NFR ID-г энэ дизайн хангах вэ?
Хэн унших вэ?{]}}

Эхлэл: Nogoolin; \href{../../requirements/scope-charter.md}{Scope
Charter}, \href{../../requirements/requirements.md}{таван шаардлага}.
Уншигч нь хөгжүүлэгч Тэнгис, хичээлийн багш.

\subsection{2. Системийн орчин ба
бүрэлдэхүүн}\label{doc6-ux441ux438ux441ux442ux435ux43cux438ux439ux43d-ux43eux440ux447ux438ux43d-ux431ux430-ux431ux4afux440ux44dux43bux434ux44dux445ux4afux4afux43d}

\emph{{[}Бөглөх: системийн хил, гаднын оролцогч, бүрэлдэхүүн бүрийн
үүрэг ба холбоо. Context diagram нэм.{]}}

Эх материал: \href{https://github.com/Tengis01/Nogoolin}{Nogoolin public
repository}, web/admin, API, өгөгдөл/auth, shared validation;
\href{../../phase-0/README.md}{Phase 0 index}. Эдгээрийг бэлэн
архитектурын нотолгоо гэж үзэхгүй, холбогдох кодтой тулгана.

\subsection{3. Архитектурын
шийдвэр}\label{doc6-ux430ux440ux445ux438ux442ux435ux43aux442ux443ux440ux44bux43d-ux448ux438ux439ux434ux432ux44dux440}

\emph{{[}Бөглөх: шийдвэр, хувилбарууд, сонгосон шалтгаан, сул тал,
хангах requirement ID.{]}}

Карт: \textbf{Decision ID:} \nolinkurl{[D-XX]}; \textbf{Requirement:}
\nolinkurl{[FR/NFR-XX]}; \textbf{Сонголт:} \nolinkurl{[…]};
\textbf{Үндэслэл/эрсдэл:} \nolinkurl{[…]}; \textbf{Source:}
\nolinkurl{[…]}.

Next.js, Fastify, Supabase зэрэг технологийн сонголтыг
\href{../../../agent-context/DECISIONS.md}{шийдвэрийн эх}-ээс нягталж
энд оруулна. ``Хэрэглэгч юу хийж чадах вэ?'' гэсэн шаардлагыг
технологийн нэрээр орлуулахгүй.

\subsection{4. Өгөгдөл ба
хамгаалалт}\label{doc6-ux4e9ux433ux4e9ux433ux434ux4e9ux43b-ux431ux430-ux445ux430ux43cux433ux430ux430ux43bux430ux43bux442}

\emph{{[}Бөглөх: entity/холбоо, ownership, эрх, өгөгдлийн lifecycle;
шаардлагын ID ба ER эх.{]}}

Эхлэх жишээ: NFR-01-ийн хувийн өгөгдөл тусгаарлах шаардлагад session
identity, ownership filter, RLS/RBAC ямар үүрэгтэйг тайлбарлана.

\subsection{5. Интерфэйс ба ажиллах
урсгал}\label{doc6-ux438ux43dux442ux435ux440ux444ux44dux439ux441-ux431ux430-ux430ux436ux438ux43bux43bux430ux445-ux443ux440ux441ux433ux430ux43b}

\emph{{[}Бөглөх: интерфэйсийн contract холбоос, бүрэлдэхүүн хоорондын
дараалал, хэвийн/алдааны урсгал.{]}}

FR-01-ийн асуулга илгээх урсгалаас эхэлнэ. API field/schema-г нэг эхэд
хадгалж, энд холбоно; \href{../../phase-0/06-api-spec.yaml}{OpenAPI}-г
кодтой тулгах шаардлагатай.

\subsection{6. Чанарын шаардлагыг хангах
арга}\label{doc6-ux447ux430ux43dux430ux440ux44bux43d-ux448ux430ux430ux440ux434ux43bux430ux433ux44bux433-ux445ux430ux43dux433ux430ux445-ux430ux440ux433ux430}

\emph{{[}Бөглөх: NFR ID → дизайны арга → шалгах арга → эрсдэл.{]}}

Жишээ чиглэл: NFR-02-д fallback, NFR-03-д тестийн хаалт. Арга төлөвлөсөн
эсэх, хэрэгжсэн эсэх, шалгасан эсэхийг ялгана.

\subsection{7. Байршуулалт ба
хамаарал}\label{doc6-ux431ux430ux439ux440ux448ux443ux443ux43bux430ux43bux442-ux431ux430-ux445ux430ux43cux430ux430ux440ux430ux43b}

\emph{{[}Бөглөх: local/test/production орчин, компонент байрших газар,
build/deploy хамаарал, тохиргооны нэр.{]}} Нууц түлхүүрийн утга
бичихгүй. \href{../../phase-0/09-deployment.md}{Deployment эх}-ийг
ашиглана.

\subsection{8. Нээлттэй асуудал ба
төлөвлөгөө}\label{doc6-ux43dux44dux44dux43bux442ux442ux44dux439-ux430ux441ux443ux443ux434ux430ux43b-ux431ux430-ux442ux4e9ux43bux4e9ux432ux43bux4e9ux433ux4e9ux4e9}

\begin{longtable}[]{@{}
  >{\raggedright\arraybackslash}p{(\linewidth - 4\tabcolsep) * \real{0.2600}}
  >{\raggedright\arraybackslash}p{(\linewidth - 4\tabcolsep) * \real{0.5500}}
  >{\raggedright\arraybackslash}p{(\linewidth - 4\tabcolsep) * \real{0.1900}}@{}}
\toprule\noalign{}
\begin{minipage}[b]{\linewidth}\raggedright
Хэсэг
\end{minipage} & \begin{minipage}[b]{\linewidth}\raggedright
Бөглөх зүйл
\end{minipage} & \begin{minipage}[b]{\linewidth}\raggedright
Status
\end{minipage} \\
\midrule\noalign{}
\endhead
\bottomrule\noalign{}
\endlastfoot
§1--4 & Архитектурын тайлбар, шийдвэр, өгөгдөл/эрхийн дизайн & planned
(W04) \\
§5 & API contract нийцүүлэлт & planned (W05) \\
§2/4/5 зураг & Context, ER, sequence диаграмын эх & planned (W08) \\
§6--7 & CI/build шалгалтын нотолгоо & planned (W09) \\
\end{longtable}

Эдгээр нь дараагийн ажлын төлөвлөгөө; M2-д дээрх бөглөх бүтэц бэлэн
болсон. Шийдвэр өөрчлөгдвөл SRS ID, traceability, холбогдох диаграмыг
хамтад нь шинэчилнэ.


\end{document}
